\documentclass[12pt]{amsart}
\usepackage{amsfonts,latexsym,rawfonts,amsmath,amssymb,amsthm,times,a4wide}
\usepackage{graphicx}
\usepackage[plainpages=false]{hyperref}
\usepackage{mathrsfs}
\usepackage{enumerate}
\usepackage{bbm}

\numberwithin{equation}{section}
\DeclareMathOperator{\arccosh}{arccosh}

\usepackage{amscd}
\usepackage{eufrak}
\usepackage{euscript}
\usepackage{array}
\usepackage{color}
\usepackage{stmaryrd}

\renewcommand\baselinestretch{1.1}
\renewcommand{\theequation}{\arabic{section}.\arabic{equation}}

\newcommand{\beq}{\begin{equation}}
\newcommand{\eeq}{\end{equation}}
\newcommand{\beqs}{\begin{eqnarray*}}
\newcommand{\eeqs}{\end{eqnarray*}}
\newcommand{\beqn}{\begin{eqnarray}}
\newcommand{\eeqn}{\end{eqnarray}}

\newtheorem{prop}{Proposition}[section]
\newtheorem{theo}[prop]{Theorem}
\newtheorem{lem}[prop]{Lemma}
\newtheorem{claim}[prop]{Claim}
\newtheorem{cor}[prop]{Corollary}
\newtheorem{rem}[prop]{Remark}
\newtheorem{ex}[prop]{Example}
\newtheorem{defi}[prop]{Definition}
\newtheorem{conj}[prop]{Conjecture}

\newcommand{\R}{\mathbb{R}}
\newcommand{\N}{\mathbb{N}}
\newcommand{\Z}{\mathbb{Z}}
\newcommand{\C}{\mathbb{C}}
\newcommand{\T}{\mathcal{T}}
\newcommand{\wt}{\widetilde}

\newcommand{\tm}{\begin{theo}}
\newcommand{\tmd}{\end{theo}}
\newcommand{\co}{\begin{cor}}
\newcommand{\cod}{\end{cor}}
\newcommand{\prp}{\begin{prop}}
\newcommand{\prpd}{\end{prop}}

\begin{document}

\title[The character of Thurston's circle packings]{The character of Thurston's circle packings}

\author{Huabin Ge}
\address{School of Mathematics, Renmin University of China, Beijing 100872, China}
\email{hbge@ruc.edu.cn}

\author{Aijin Lin}
\address{Department of Mathematics, National University of Defense Technology, Changsha 410073, China}
\email{linaijin@nudt.edu.cn}

\begin{abstract}
We introduce the character of Thurston's circle packings in hyperbolic geometry background.
Consequently, some quite simple criteria are obtained for the existence of hyperbolic circle packings. For example,
if a closed surface $X$ admits a circle packing with all vertex degree $d_i \geq 7$, then it admits a unique complete
hyperbolic metric, so that the triangulation graph of the circle packing is isotopic to a geometric decomposition
of $X$. This criterion is sharp due to the fact that any closed hyperbolic surface admits no triangulations with
all $d_i \leq 6$. As a corollary, we obtain a new proof of the uniformization theorem for closed surfaces with genus
$g \geq 2$, and moreover, any hyperbolic closed surface has a geometric decomposition. To approach our results, we
shall use Chow-Luo's combinatorial Ricci flow as a fundamental tool.
\end{abstract}

\keywords{circle packings, combinatorial Ricci flow, character}

\maketitle

\section{Background}

Thurston's hyperbolization theorem for 3-manifolds is one of the truly great mathematical discoveries of
the twentieth century. It establishes a deep and strong link between the geometry and topology of 3-
manifolds and the algebra of discrete groups of $\mathrm{Isom}(\mathbb{H}^3)$. Proving Thurston's hyperbolization for Haken
3-manifolds requires three major tools: the existence of hierarchies in Haken manifolds which allows us
to cut a Haken manifold into polyhedra; Andreev's theorem which allows us to give these polyhedra
hyperbolic structures; and the skinning lemma which allows us to glue the pieces together again.

The beautiful theorem of Andreev \cite{Andreev1,Andreev2} states that if we make a topological model for a polyhedron
and choose candidate dihedral angles (angles between the faces) that are at most $\pi/2$, then there are
simply verifiable conditions that tell us whether there exists a hyperbolic polyhedron with the assigned
angles. Furthermore, if such a polyhedron exists, it is unique. Andreev's theorem provides a complete
characterization of compact hyperbolic polyhedra with non-obtuse dihedral angles, and is essential for
proving Thurston's hyperbolization theorem for Haken 3-manifolds.

Thurston \cite{Thurston} observed a very deep connection between circle packings and hyperbolic polyhedra. Given
a convex hyperbolic polyhedron in the hyperbolic 3-space $\mathbb{B}^3$, the boundaries of the oriented hyperbolic
planes containing its faces form a circle packing on the sphere $\partial \mathbb{B}^3$. This circle packing records all the
information of the original polyhedron: its combinatorial type is exactly dual to the 1-skeleton of the
polyhedron, and the exterior intersection angles between the circles are equal to the dihedral angles
between the faces of the hyperbolic polyhedron. Thurston's circle packing theorem regards the existence
and uniqueness of circle packings on higher genus surfaces with a prescribed combinatorial type and non-
obtuse exterior intersection angles. Thurston pointed out that the sphere version of his theorem followed
from the Andreev's theorem. For related results on circle packings and hyperbolic polyhedra, we refer to
the works of Marden-Rodin \cite{MardenRodin}, Colin de Verdi\`ere \cite{Colin}, Bowditch \cite{Bowditch}, Rivin-Hodgson \cite{RivinHodgson}, Rivin \cite{Rivin1,Rivin2},
Bao-Bonahon \cite{BaoBonahon}, Bobenko-Springborn \cite{BobenkoSpringborn}, Leibon \cite{Leibon}, Rousset \cite{Rousset}, Schlenker \cite{Schlenker}
and others. Thurston also posed a conjecture regarding the convergence of infinitesimal hexagonal tangent circle packings to
conformal mappings \cite{Thurston2}, which was proved by Rodin-Sullivan \cite{RodinSullivan}. From then on circle packings have
played significant roles in the study of low dimensional geometry and topology, complex analysis (such as
the famous Koebe uniformization conjecture for circle domains, see e.g. \cite{HeSchramm1,HeSchramm2,HeSchramm3}), and various problems
in combinatorics \cite{LiuZhou,Schramm1}, discrete and computational geometry \cite{DaiGuLuo,Stephenson}, minimal surfaces \cite{BobenkoHoffmannSpringborn}, and many
others.

A circle packing $\mathcal{P} = \{C_v : v \in V\}$ on a surface is a collection of circles with a particular combinatorial
structure. Let $X$ be a closed surface with a triangulation $\mathcal{T} = (V,E,F)$, where $V,E,F$ are the set
of vertices, edges and triangles respectively. Assume $\mu$ is a Riemannian metric on $X$ with constant
curvature. Given a weight $\Phi : E \to [0,\pi/2]$, a circle packing $\mathcal{P}$ on $(X,\mu)$ is called $(\mathcal{T},\Phi)$-type if there
exists a geodesic triangulation of $\mathcal{T}$ on $(X,\mu)$ isotopic to $\mathcal{T}$ such that the circle $C_v$ is centered at $\mathcal{T}_\mu(v)$
and for any edge $e \in E$, the two circles $C_v, C_u$, which correspond to the vertices $v,u$ of $e$, intersect at
an angle $\Phi(e)$. A fundamental question arises naturally: does there exist a $(\mathcal{T},\Phi)$-type circle packing $\mathcal{P}$
on $(X,\mu)$, and is it unique when it exists? A celebrated answer to this question is the following circle
packing theorem due to Thurston \cite[Chap.~13, Theorem~13.7.1]{Thurston}.

\begin{theo}[Thurston]\label{thm:Thurston}
Let $\mathcal{T}$ be a triangulation of a closed surface $X$ of genus $g > 0$. Let
$\Phi: E \to [0,\pi/2]$ be a function satisfying the following conditions:
\begin{itemize}
\item[($\mathrm{T}_1$).] If $e_1, e_2, e_3$ form a null-homotopic closed path in $\mathcal{T}$, and if $\sum_{l=1}^3 \Phi(e_l) \geq \pi$, then these edges form
the boundary of a triangle of $\mathcal{T}$;
\item[($\mathrm{T}_2$).] If $e_1, e_2, e_3, e_4$ form a null-homotopic closed path and if
$\sum_{l=1}^4 \Phi(e_l) = 2\pi$, then $e_1, e_2, e_3, e_4$ form the
boundary of the union of two adjacent triangles.
\end{itemize}
Then there is a constant curvature metric $\mu$ on $X$ such that $(X,\mu)$ supports a $(\mathcal{T},\Phi)$-type circle packing
$\mathcal{P}$. Moreover, the pair $(\mu,\mathcal{P})$ is unique up to isometries if $g > 1$, and up to similarities if $g = 1$.
\end{theo}

The sphere version of the above theorem follows from Koebe \cite{Koebe} and Andreev \cite{Andreev1,Andreev2}. Together they are
often named Koebe-Andreev-Thurston's circle packing theorem. In this paper we call the above theorem
Thurston's circle packing theorem for simplicity. Of course this will not weaken Koebe and Andreev's
contributions. Thurston's theorem gives a complete criterion of the existence of circle packings. Note
it does not assume the existence of the metric $\mu$ a priori. Consequently, there is a uniquely determined
geometric decomposition of $(X,\mu)$ isotopic to $\mathcal{T}$ so that the edges are geodesics.

\section{Main Results}

Theorem \ref{thm:Thurston} gives a wonderful criterion for the existence of circle packings with constant curvatures.
However, Thurston's criteria $(\mathrm{T}_1)$ and $(\mathrm{T}_2)$ are extremely difficult to verify for general weight $\Phi \in [0,\pi/2]$,
since the angle structure $\Phi$ and the combinatorial structure of $\mathcal{T}$ are globally intertwined together on $X$.
Inspired by the work in Ge-Jiang-Shen \cite{GeJiangShen}, Ge-Hua \cite{GeHua} and Feng-Ge-Hua \cite{FengGeHua} on 3-dimensional geometric
triangulations, we introduce the characters of circle packings to overcome this difficulty. Consequently,
we found some quite simple criteria for the existence of Thurston's circle packings. First we have

\begin{theo}\label{thm:main1}
Let $X$ be a closed surface. If $X$ admits a triangulation $\mathcal{T}$ with degree $d_i \geq 7$ at each
vertex, then for any given constant weight $\Phi : E \to [0,\pi/2]$ there exists a unique complete hyperbolic
metric $\mu$ on $X$ so that $(X,\mu)$ supports a $(\mathcal{T},\Phi)$-type circle packing $\mathcal{P}$.
\end{theo}

The requirement that $d_i \geq 7$ is sharp in the sense that every circle packing $\mathcal{P}$ with all vertex degree
$d_i \leq 6$ can not be supported on any closed surface $X$ of genus $g \geq 2$, and every circle packing $\mathcal{P}$ with all
vertex degree $d_i \geq 7$ can not be supported on the sphere or torus, see Section \ref{sec:circle} for details.

The requirement that $\Phi$ is a constant in Theorem \ref{thm:main1} is too restrictive, and it can be released to some
extent. In fact, Theorem \ref{thm:main1} can be generalized to the following

\begin{theo}\label{thm:main2}
Let $X$ be a closed surface. If $X$ admits a triangulation $\mathcal{T}$ with degree $d_i \geq 7$ at each
vertex, then for any weight $\Phi \in [\arccos\eta, \arccos\xi] \subset [0,\pi/2]$, where $0 \leq \xi \leq \eta \leq 1$ are arbitrary chosen
so that $\eta < (2\cos\frac{2\pi}{7} - 1 + \xi)/(2 - 2\cos\frac{2\pi}{7})$, there exists a unique complete hyperbolic metric $\mu$ on $X$ so
that $(X,\mu)$ supports a $(\mathcal{T},\Phi)$-type circle packing $\mathcal{P}$.
\end{theo}

Assume $\Phi$ is a constant, or $\Phi$ takes values in $[0, 0.33\pi]$ or in $[0.4\pi, \pi/2]$ respectively. They all satisfy the
assumption in Theorem \ref{thm:main2} by Corollary \ref{cor:3.8}. However, under a slightly stronger condition than $d_i \geq 7$,
we do not need any restrictions on $\Phi$.

\begin{theo}\label{thm:main3}
Let $X$ be a closed surface. If $X$ admits a triangulation $\mathcal{T}$ with degree $d_i \geq 9$ at each
vertex, then for any given weight $\Phi: E \to [0,\pi/2]$ there exists a unique complete hyperbolic metric $\mu$ on
$X$ such that $(X,\mu)$ supports a $(\mathcal{T},\Phi)$-type circle packing $\mathcal{P}$.
\end{theo}

Compared to the Thurston's conditions $(\mathrm{T}_1)$ and $(\mathrm{T}_2)$, our degree conditions do not respect the angle
structure $\Phi$. This seems quite amazing and may lie deep. Actually, the two degree criteria (i.e. $d_i \geq 7$ for
constant weight or $d_i \geq 9$ for arbitrary weight) are all special cases of the ``character criteria'' $L_i(\mathcal{T},\Phi) >
2\pi$ for each vertex $i \in V$, where
\begin{equation}\label{eq:character}
L_i(\mathcal{T},\Phi) = \sum_{\triangle ijk \in F} \arccos \frac{1 + \cos\Phi_{ij} + \cos\Phi_{ik} - \cos\Phi_{jk}}{2\sqrt{1 + \cos\Phi_{ij}}\sqrt{1 + \cos\Phi_{ik}}},
\end{equation}
where the sum runs over all triangles having a vertex $i$. The character of $(X,\mathcal{T},\Phi)$ is
\[
L(\mathcal{T},\Phi) = (L_1(\mathcal{T},\Phi), \cdots, L_N(\mathcal{T},\Phi)).
\]
We will show in Section \ref{sec:circle} that Theorem \ref{thm:main1},
\ref{thm:main2} and Theorem \ref{thm:main3} are all special cases of the following theorem.

\begin{theo}\label{thm:main4}
Let $X$ be a closed surface with a triangulation $\mathcal{T}$ and a weight $\Phi : E \to [0,\pi/2]$. If
the character $L_i(\mathcal{T},\Phi) > 2\pi$ at all vertices, then $\chi(X) < 0$ and there exists a unique complete hyperbolic
metric $\mu$ on $X$ such that $(X,\mu)$ supports a $(\mathcal{T},\Phi)$-type circle packing $\mathcal{P}$. On the other hand, if the
character $L_i(\mathcal{T},\Phi) \leq 2\pi$ at all vertices, then there exists no such $(\mathcal{T},\Phi)$-type circle packings.
\end{theo}

The distinct feature of the characters $L_i(\mathcal{T},\Phi)$ is their localities. They can be verified locally and
separately at each vertex. However, Thurston's conditions $(\mathrm{T}_1)$ and $(\mathrm{T}_2)$ are intertwined together, which
are global and can not be verified separately.

As an application of Theorem \ref{thm:main1}, we reprove the following well known uniformization theorem.
Actually, combining Theorem \ref{thm:main1} in the case of weight $\Phi \equiv 0$ with the fact that every orientable
closed surface of genus $g \geq 2$ admits a triangulation with degree $d = 7$ at each vertex, we obtain
the uniformization theorem for surfaces with higher genus:

\begin{cor}\label{cor:uniformization}
Every closed surface of genus $g \geq 2$ admits a complete hyperbolic metric.
\end{cor}

Moreover, we reprove the following geometric decomposition theorem for hyperbolic surfaces. We refer
to Martelli's book \cite{Martelli} for more details on geometric decompositions.

\begin{cor}\label{cor:geometric}
A closed surface $X$ of genus $g \geq 2$ admits infinitely many geometric decompositions.
Actually, for any triangulation $\mathcal{T}$ with degree $d > 6$ at each vertex, there exists a geometric decomposition
on $X$ isotopic to $\mathcal{T}$ so that the edges are geodesics.
\end{cor}

\section{Circle packings and their characters}\label{sec:circle}

Circle packings provide a bridge between the combinatoric on one hand, and the geometry on the other
hand. Let $X$ be a closed surface with a triangulation $\mathcal{T} = (V,E,F)$, where $V,E,F$ denote the sets
of vertices, edges, and faces respectively. A weight on the triangulation is defined to be a function
$\Phi : E \to [0,\pi/2]$. Throughout this paper, a function defined on vertices is an $N$-dimensional column
vector, where $N = |V|$ is the number of vertices. Moreover, all vertices $v_1, \ldots, v_N$ are abbreviated as
$1, \ldots, N$ for simplicity. Next we study the circle packings on a topological surface from the perspective
of metrics.

Endow $X$ with hyperbolic (flat, resp.) cone metrics as follows. A circle packing metric based on
$(X,\mathcal{T},\Phi)$ is a function $r : V \to (0,+\infty)$. We consider $r = (r_1, \cdots, r_N) \in \mathbb{R}^N_{>0}$ as a point in $\mathbb{R}^N_{>0} =
(0,+\infty)^N$ in the whole paper. Geometrically, a circle packing metric $r$ assigns a circle with radius $r_i$ at
every vertex $i$, and realizes each edge $e_{ij} \in E$ joining $i$ to $j$ by a hyperbolic (Euclidean, resp.) segment
of length
\begin{equation}\label{eq:hyp-edge}
l_{ij} = \cosh^{-1}(\cosh r_i \cosh r_j + \sinh r_i \sinh r_j \cos\Phi(e_{ij}))
\end{equation}
(resp. $l_{ij} = \sqrt{r_i^2 + r_j^2 + 2r_i r_j \cos\Phi(e_{ij})}$).

Thurston once noted \cite[Lemma~13.7.2]{Thurston} that, for each triangle $\triangle ijk \in F$, the three edge lengths $l_{ij}$,
$l_{jk}$ and $l_{ik}$ satisfy triangle inequalities. Thus one can realize each triangle $\triangle ijk \in F$ by a hyperbolic
(Euclidean, resp.) triangle of edge lengths $l_{ij}, l_{jk}$ and $l_{ik}$. The triangle is formed by the centers of three
circles of radii $r_i, r_j$ and $r_k$ intersecting at angles $\Phi_{ij}, \Phi_{jk}$ and $\Phi_{ik}$, where $\Phi_{mn}$ denotes the value of $\Phi$
on the edge $e_{mn} \in E$. This produces a hyperbolic cone metric (flat cone metric, resp.) on the surface
$X$ with singularities at the vertices. Denote $\theta_i^{jk}$ (or $\theta_i$ when there is no confusion) as the inner angle at
a vertex $i$ in the triangle $\triangle ijk \in F$. Let $A_i = \sum_{\triangle ijk \in F} \theta_i^{jk}$ be the cone angle at the vertex $i$, which is
equal to the sum of inner angles at $i$ for all triangles incident to $i$. The combinatorial Gauss curvature
$K_i$ at $i$ is defined as
\begin{equation}\label{eq:curvature}
K_i = 2\pi - A_i = 2\pi - \sum_{\triangle ijk \in F} \theta_i^{jk}.
\end{equation}
Then for a given $(\mathcal{T},\Phi)$-type circle packing $\mathcal{P}$, every circle packing metric $r = (r_1, \cdots, r_N)$ produces a
hyperbolic (flat, resp.) cone metric on $X$ with cone singularities $A_i$ centered on each $i \in V$. There are no
singularities on $X \setminus V$ by the constructions. At a particular vertex $i$, $A_i = 2\pi$ means there is no singularity
at $i$. Hence a circle packing metric $r_{ze}$ with $K(r_{ze}) = 0$ is particular meaningful in the sense that there
are no singularities on $X$, even on $V$: it produces a complete hyperbolic (flat, resp.) metric $\mu$ on $X$. It
is $\mathcal{T}$-type and has exterior intersection angles given by $\Phi$. Moreover, those hyperbolic (Euclidean, resp.)
triangles form a geodesic triangulation of $(X,\mu)$. Specifically, when it produces a hyperbolic metric, we
name the unique circle packing $r_{ze}$ a hyperbolic circle packing.

\begin{rem}
By abuse of language, a circle packing metric $r$ is also called a circle packing for brevity.
The original $(\mathcal{T},\Phi)$-type circle packing $\mathcal{P}$ may be considered as the equivalent class of all circle packing
metrics $r \in \mathbb{R}^N_{>0}$. We say our setting is in the hyperbolic background if we consider hyperbolic cone
metrics constructed by gluing hyperbolic triangles with hyperbolic edge lengthes \eqref{eq:hyp-edge}. The Euclidean
background is defined similarly.
\end{rem}

\begin{defi}\label{def:character}
Let $X$ be a closed surface with a triangulation $\mathcal{T}$ and a weight $\Phi \in [0,\pi/2]$. For any
circle packing based on $(X,\mathcal{T},\Phi)$, define the character $L_i(\mathcal{T},\Phi)$ at each $i \in V$ by
\begin{equation}\label{eq:char-def}
L_i(\mathcal{T},\Phi) = \sum_{\triangle ijk \in F} \arccos \frac{1 + \cos\Phi_{ij} + \cos\Phi_{ik} - \cos\Phi_{jk}}{2\sqrt{1 + \cos\Phi_{ij}}\sqrt{1 + \cos\Phi_{ik}}},
\end{equation}
where the sum runs over all triangles having a vertex $i$. The character of $(X,\mathcal{T},\Phi)$ is
\[
L(\mathcal{T},\Phi) = (L_1(\mathcal{T},\Phi), \cdots, L_N(\mathcal{T},\Phi)).
\]
\end{defi}

It seems that the character is just defined for a weighted triangulation $(\mathcal{T},\Phi)$ on $X$, and has no relation
to any circle packings. However, both the combinatorial structure $\mathcal{T}$ and the angle structure $\Phi$ come
from a circle packing $\mathcal{P}$ on $X$. Recall that a circle packing $\mathcal{P} = \{C_v : v \in V\}$ on $(S,\mu)$ is of $(\mathcal{T},\Phi)$-type if
there exists a geodesic triangulation of $\mathcal{T}$ on $(S,\mu)$ isotopic to $\mathcal{T}$ such that the circle $C_v$ is centered at $\mathcal{T}_\mu(v)$
and for any edge $e \in E$, the two circles $C_v, C_u$, which correspond to the vertices $v,u$ of $e$, intersect at an
angle $\Phi(e)$. Hence the weighted triangulation $(\mathcal{T},\Phi)$ records all the information of a $(\mathcal{T},\Phi)$-type circle
packing $\mathcal{P}$, that is, the combinatorial structure given by $\mathcal{T}$ and the angle structure given by $\Phi$. Hence
the character $(L_i(\mathcal{T},\Phi))_{i \in V}$ is indeed an invariant of all $(\mathcal{T},\Phi)$-type circle packings $\mathcal{P}$ on a closed surface
$X$.

\begin{prop}\label{prop:3.3}
Given a weighted triangulated surface $(X,\mathcal{T},\Phi)$. For each vertex $i \in V$, the
character $L_i(\mathcal{T},\Phi)$ is exactly the cone angle $A_i$ of a circle packing metric $r = (1,\cdots, 1)$ in the Euclidean
background.
\end{prop}

\begin{proof}
Assume $r_i = 1$ for each vertex $i \in V$. Then by the Euclidean edge length formula, $l_{ij} = \sqrt{2 + 2\cos\Phi_{ij}}$ for each edge $e_{ij} \in E$.
Consider a triangle $\triangle ijk \in F$, it's easy to see
\[
\cos\theta_i^{jk} = \frac{1 + \cos\Phi_{ij} + \cos\Phi_{ik} - \cos\Phi_{jk}}{2\sqrt{1 + \cos\Phi_{ij}}\sqrt{1 + \cos\Phi_{ik}}}.
\]
Then by the definition of the character \eqref{eq:char-def}, we obtain $L_i(\mathcal{T},\Phi) = \sum_{\triangle ijk \in F} \theta_i^{jk} = A_i$.
\end{proof}

\begin{prop}\label{prop:3.4}
Denote the average character by $L_{av} = \sum_{i \in V} L_i(\mathcal{T},\Phi) / N$, then
\[
L_{av} = 2\pi\left(1 - \frac{\chi(X)}{N}\right).
\]
Consequently, $L_{av} > 2\pi$ if $\chi(X) < 0$, $L_{av} = 2\pi$ if $\chi(X) = 0$ and $L_{av} < 2\pi$ if $\chi(X) > 0$.
\end{prop}

\begin{proof}
Note $2\pi - L_i(\mathcal{T},\Phi)$ is the combinatorial curvature $K_i$ of a particular packing $r \equiv 1$ in the
Euclidean background by Proposition \ref{prop:3.3}. The conclusion follows from the following combinatorial Gauss-
Bonnet formula \cite{ChowLuo}
\[
\sum_{i \in V} K_i = 2\pi\chi(X)
\]
in the Euclidean background. We can also prove this proposition directly. For simplicity, we write
\[
\gamma_i^{jk} = \arccos \frac{1 + \cos\Phi_{ij} + \cos\Phi_{ik} - \cos\Phi_{jk}}{2\sqrt{1 + \cos\Phi_{ij}}\sqrt{1 + \cos\Phi_{ik}}}.
\]
Consider a particular triangle $\triangle ijk \in F$ with edge lengths $l_{mn} = \sqrt{2 + 2\cos\Phi_{mn}}$, where an edge $e_{mn}$ is
chosen from $\{e_{ij}, e_{jk}, e_{ki}\}$. Note $\gamma_i^{jk}$ is just the inner angle at the vertex $i$. Hence we have
\[
\gamma_i^{jk} + \gamma_j^{ik} + \gamma_k^{ij} = \pi.
\]
Since $X$ is closed and $\mathcal{T}$ is a triangulation, then $2|E| = 3|F|$, and it follows
\begin{align*}
\sum_{i \in V} L_i(\mathcal{T},\Phi) &= \sum_{i \in V} \sum_{\triangle ijk \in F} \gamma_i^{jk} \\
&= \sum_{\triangle ijk \in F} (\gamma_i^{jk} + \gamma_j^{ik} + \gamma_k^{ij}) \\
&= \pi |F| = 2\pi(|E| - |F|) = 2\pi N - 2\pi\chi(X).
\end{align*}
\end{proof}

Obviously, Theorem \ref{thm:main1} follows from Theorem \ref{thm:main4} and the following corollary:

\begin{cor}\label{cor:3.5}
Given a weighted triangulated surface $(X,\mathcal{T},\Phi)$ such that $\Phi : E \to [0,\pi/2]$ is a
constant. If $d_i \geq 7$ at a vertex $i \in V$, then $L_i(\mathcal{T},\Phi) > 2\pi$.
\end{cor}

\begin{proof}
If $\Phi$ takes a constant, then $L_i(\mathcal{T},\Phi) = \sum_{\triangle ijk \in F} \arccos \frac{1}{2} = \frac{\pi d_i}{3} \geq \frac{7\pi}{3}$ by \eqref{eq:char-def}.
\end{proof}

\begin{prop}\label{prop:3.6}
Given a weighted triangulated surface $(X,\mathcal{T},\Phi)$ with $\Phi : E \to [0,\pi/2]$. For each
vertex $i \in V$, the character has a lower bound
\[
L_i(\mathcal{T},\Phi) \geq d_i \arccos \frac{3}{4} \approx \frac{2.07\pi}{9} d_i.
\]
\end{prop}

\begin{proof}
For arbitrary $a, b \in [0,1]$, it is easy to show $4a^2 + 4b^2 \leq 5 + a + b + ab$. It follows
\[
\frac{(1 + a + b)^2}{(1 + a)(1 + b)} \leq \frac{9}{4}.
\]
Now consider a triangle $\triangle ijk \in F$ and substitute $a = \cos\Phi_{ij}$ and $b = \cos\Phi_{ik}$ into the above inequality,
we get
\[
\left(\frac{1 + \cos\Phi_{ij} + \cos\Phi_{ik} - \cos\Phi_{jk}}{2\sqrt{1 + \cos\Phi_{ij}}\sqrt{1 + \cos\Phi_{ik}}}\right)^2 \leq \frac{(1 + \cos\Phi_{ij} + \cos\Phi_{ik})^2}{4(1 + \cos\Phi_{ij})(1 + \cos\Phi_{ik})} \leq \frac{9}{16},
\]
which implies
\begin{equation}\label{eq:3.5}
0 \leq \frac{1 + \cos\Phi_{ij} + \cos\Phi_{ik} - \cos\Phi_{jk}}{2\sqrt{1 + \cos\Phi_{ij}}\sqrt{1 + \cos\Phi_{ik}}} \leq \frac{3}{4}.
\end{equation}
It follows the conclusion.
\end{proof}

Theorem \ref{thm:main3} follows from Theorem \ref{thm:main4} and the following corollary:

\begin{cor}\label{cor:3.7}
Given a weighted triangulated surface $(X,\mathcal{T},\Phi)$ with $\Phi: E \to [0,\pi/2]$. If $d_i \geq 9$ at a
vertex $i \in V$, then $L_i(\mathcal{T},\Phi) > 2\pi$.
\end{cor}

By thorough analysis on the proof of Proposition \ref{prop:3.6}, Corollary \ref{cor:3.5} can be improved. In fact, Theorem
\ref{thm:main2} follows from Theorem \ref{thm:main4} and the following corollary:

\begin{cor}\label{cor:3.8}
Given a weighted triangulated surface $(X,\mathcal{T},\Phi)$. Assume the weight satisfies $\Phi \in
[0, 0.33\pi]$, or $\Phi \in [0.4\pi, \pi/2]$, or more generally $\Phi \in [\arccos\eta, \arccos\xi] \subset [0,\pi/2]$, where $0 \leq \xi \leq \eta \leq 1$
are arbitrary chosen so that $\eta < (2\cos\frac{2\pi}{7} - 1 + \xi)/(2 - 2\cos\frac{2\pi}{7})$. If $d_i \geq 7$ at a vertex $i \in V$, then
$L_i(\mathcal{T},\Phi) > 2\pi$.
\end{cor}

\begin{proof}
Denote $\lambda = 2\cos\frac{2\pi}{7}$. For arbitrary $a, b \in [\xi, \eta] \subset [0,1]$, it is easy to prove
\[
1 + a + b - \lambda\sqrt{1 + a}\sqrt{1 + b} \leq 1 + 2\eta - \lambda(1 + \eta).
\]
Further assume $\eta < \frac{\lambda - 1 + \xi}{2 - \lambda}$, which is equivalent to $1 + 2\eta - \lambda(1 + \eta) < \xi$, then
\[
1 + a + b - \lambda\sqrt{1 + a}\sqrt{1 + b} < \xi.
\]
Assume $\Phi \in [\arccos\eta, \arccos\xi] \subset [0,\pi/2]$, then $\xi \leq \cos\Phi_{ij}, \cos\Phi_{ik}, \cos\Phi_{jk} \leq \eta$. Substitute $a = \cos\Phi_{ij}$
and $b = \cos\Phi_{ik}$ into the above inequality, we get
\[
1 + \cos\Phi_{ij} + \cos\Phi_{ik} - \lambda\sqrt{1 + \cos\Phi_{ij}}\sqrt{1 + \cos\Phi_{ik}} < \xi \leq \cos\Phi_{jk}.
\]
Then it follows
\[
0 \leq \frac{1 + \cos\Phi_{ij} + \cos\Phi_{ik} - \cos\Phi_{jk}}{2\sqrt{1 + \cos\Phi_{ij}}\sqrt{1 + \cos\Phi_{ik}}} < \frac{\lambda}{2} = \cos\frac{2\pi}{7}
\]
and further
\[
\arccos \frac{1 + \cos\Phi_{ij} + \cos\Phi_{ik} - \cos\Phi_{jk}}{2\sqrt{1 + \cos\Phi_{ij}}\sqrt{1 + \cos\Phi_{ik}}} > \frac{2\pi}{7}.
\]
If $d_i \geq 7$, then we get the conclusion
\[
L_i(\mathcal{T},\Phi) > \sum_{\triangle ijk \in F} \frac{2\pi}{7} \geq 7 \cdot \frac{2\pi}{7} = 2\pi.
\]
Specially, choose $[\xi, \eta] = [0, 0.309]$ or $[\xi, \eta] = [0.509, 1]$. Direct calculations show that they all satisfy
the assumption $\eta < \frac{\lambda - 1 + \xi}{2 - \lambda}$ and $[\arccos 1, \arccos 0.509] = [0, 0.33\pi]$, $[\arccos 0.309, \arccos 0] = [0.4\pi, \pi/2]$.
Hence we finish the proof.
\end{proof}

\section{Combinatorial Ricci flows}\label{sec:ricci}

The combinatorial Ricci flows were introduced by Chow-Luo \cite{ChowLuo}, which are the analogue of Hamilton's
Ricci flow in the combinatorial setting. They obtained a new proof of the existence part of Thurston's
circle packing theorem, by showing that the combinatorial Ricci flows produce solutions which converge
exponentially fast to Thurston's circle packings on surfaces. Since then, the combinatorial Ricci flows
provide an effective algorithm for finding complete hyperbolic metrics, and are one of the main tools of
finding geometric structures on surfaces and 3-manifolds. Luo \cite{Luo} also initialed a program to hyperbolize
3-manifolds with boundary by the combinatorial Ricci flows. The program was formulated more clearly
and carried forward extensively by the series work of Ge and his collaborators. Particularly for a 3-
manifold with boundary of higher genus, we refer to \cite{FengGeHua} where a hyperbolic metric and an ideal geometric
decomposition were obtained under suitable combinatorial conditions.

We recall Chow-Luo's theory briefly. Let $X$ be a closed surface with a triangulation $\mathcal{T}$ and a weight
$\Phi \in [0,\pi/2]$. Considered as a smooth family of circle packings $r(t) \subset \mathbb{R}^N_{>0}$ based on $(X,\mathcal{T},\Phi)$, which evolves
according to the following combinatorial Ricci flow
\begin{equation}\label{eq:euclidean-flow}
\frac{dr_i(t)}{dt} = -K_i r_i, \quad i \in V
\end{equation}
in the Euclidean background. For any initial circle packing $r(0) \in \mathbb{R}^N_{>0}$, a solution to the combinatorial
Ricci flow \eqref{eq:euclidean-flow} is called convergent if $\lim_{t \to \infty} r_i(t) = r_i(\infty) \in \R$ exists for all $i \in V$. As a consequence,
$\lim_{t \to \infty} K_i(t) = K_i(\infty)$ exists for all $i$ since the curvature map $K(r) = (K_1(r), \cdots, K_N(r))$ is a smooth map
of $r$. A convergent solution $r(t)$ is called convergent exponentially fast if there are positive constants $c_1, c_2$
so that for all time $t \geq 0$ and all vertices $i \in V$, $|r_i(t) - r_i(\infty)| \leq c_1 e^{-c_2 t}$ and $|K_i(t) - K_i(\infty)| \leq c_1 e^{-c_2 t}$.

Chow-Luo proved that, in the Euclidean background, the solution to the normalized combinatorial Ricci
flow
\begin{equation}\label{eq:normalized-flow}
\frac{dr_i}{dt} = \left(\frac{2\pi\chi(X)}{N} - K_i\right) r_i, \quad i \in V
\end{equation}
which is related to \eqref{eq:euclidean-flow} by a change of scales, exists for all time. Moreover, the solution converges if and
only if there exists a circle packing with constant curvature $2\pi\chi(X)/N$, and in this case, the solution
converges exponentially fast to a constant curvature circle packing (which is unique up to scaling).

We will use the following combinatorial Ricci flow in the hyperbolic background
\begin{equation}\label{eq:hyperbolic-flow}
\frac{dr_i(t)}{dt} = -K_i \sinh r_i, \quad i \in V
\end{equation}
as the main tool to approach our results in this paper. It provides a useful tool to deform any initial
circle packing $r(0)$ to a hyperbolic one. In deed, in case $\chi(X) < 0$, Chow-Luo proved that the solution
to \eqref{eq:hyperbolic-flow} exists for all time $t \geq 0$ and converges if and only if Thurston's conditions $(\mathrm{T}_1)$ and $(\mathrm{T}_2)$ are
satisfied. Furthermore, if it converges, then it converges exponentially fast to a circle packing $r_{ze}$ with
$K_i = 0$ at all vertices $i \in V$, i.e. a hyperbolic circle packing. In this case, the hyperbolic circle packing
$r_{ze}$ is unique by the rigidity part of Thurston's circle packing theorem.

Thurston's criteria $(\mathrm{T}_1)$ and $(\mathrm{T}_2)$ are globally intertwined together. It is generally not easy to verify
directly. Using the character $L_i(\mathcal{T},\Phi)$ introduced in Definition \ref{def:character}, we give some simple criteria for the
existence of Thurston's hyperbolic circle packings next.

\begin{theo}\label{thm:4.1}
Given a weighted triangulated closed surface $(X,\mathcal{T},\Phi)$ with the weight $\Phi: E \to [0,\pi/2]$.
Consider the hyperbolic background setting.

(a) If the character $L_i(\mathcal{T},\Phi) > 2\pi$ at all vertices, then $\chi(X) < 0$ and there exists a unique hyperbolic
circle packing $r_{ze}$ based on $(X,\mathcal{T},\Phi)$. Moreover, the solution $r(t)$ to the combinatorial Ricci flow \eqref{eq:hyperbolic-flow}
converges exponentially fast to $r_{ze}$ for any initial circle packing $r(0) \in \mathbb{R}^N_{>0}$. Consequently, $r_{ze}$ determines
a unique complete hyperbolic metric $\mu$ on $X$ such that $(X,\mu)$ supports a geometric decomposition isotopic
to $\mathcal{T}$, and each edge connecting two adjacent vertices is a hyperbolic geodesic.

(b) If the character $L_i(\mathcal{T},\Phi) \leq 2\pi$ at all vertices, then $\chi(X) \geq 0$. Consequently, there exists no
hyperbolic circle packings based on $(X,\mathcal{T},\Phi)$. In this case, any solution $r(t)$ to \eqref{eq:hyperbolic-flow} satisfies $r(t) \to 0$
when $t$ tends to $+\infty$.
\end{theo}

\section{Geometry of hyperbolic triangles}\label{sec:triangles}

In this section, we study the geometry of hyperbolic triangles which are the basic building blocks of the
weighted triangulated surface $(X,\mathcal{T},\Phi)$. For any $\{v_i, v_j, v_k\} \subset \mathbb{H}^2$, we denote by $\triangle ijk$ the corresponding
hyperbolic triangle in $\mathbb{H}^2$. Obviously, the interior angle $\theta_i = \theta_i(\vec{r})$ at each vertex $i$ is a smooth function
of $\vec{r} = (r_i, r_j, r_k)$. By the hyperbolic cosine law we have
\[
\cos\theta_i(\vec{r}) = \frac{\cosh l_{ij} \cosh l_{ik} - \cosh l_{jk}}{\sinh l_{ij} \sinh l_{ik}}.
\]
Recall Thurston once derived
\begin{lem}[Lemma~13.7.3, \cite{Thurston}]\label{lem:Thurston}
For a weighted triangulated closed surface $(X,\mathcal{T},\Phi)$ whose weight
satisfies $\Phi : E \to [0,\pi/2]$. In the hyperbolic background geometry $\mathbb{H}^2$ one has $\partial\theta_i/\partial r_i < 0$, $\partial\theta_i/\partial r_j > 0$
for $i \neq j$ and $\partial(\theta_i + \theta_j + \theta_k)/\partial r_i < 0$.
\end{lem}

Choose a special radius vector $\vec{r} = t\vec{1} = (t, t, t)$, we get the following monotonicity proposition for the
angle function $\theta_i(t\vec{1})$.

\begin{prop}\label{prop:5.2}
The angle function $t \mapsto \theta_i(t\vec{1})$ is continuously differentiable and strictly decreasing
function in $(0, +\infty)$.
\end{prop}

\begin{proof}
Set $f(t) = \cos\theta_i(t\vec{1})$, then we have
\begin{align*}
\cosh l_{ij} &= \cosh^2 t + \sinh^2 t \cos\Phi_{ij} = \sinh^2 t(1 + \cos\Phi_{ij}) + 1; \\
\cosh l_{ik} &= \cosh^2 t + \sinh^2 t \cos\Phi_{ik} = \sinh^2 t(1 + \cos\Phi_{ik}) + 1; \\
\cosh l_{jk} &= \cosh^2 t + \sinh^2 t \cos\Phi_{jk} = \sinh^2 t(1 + \cos\Phi_{jk}) + 1.
\end{align*}
By a tedious but direct computation we have
\begin{equation}\label{eq:5.1}
f(t) = \frac{(1 + \cos\Phi_{ij})(1 + \cos\Phi_{ik})\sinh^2 t + (1 + \cos\Phi_{ij} + \cos\Phi_{ik} - \cos\Phi_{jk})}{f_1(t)f_2(t)},
\end{equation}
where
\[
f_1(t) = \sqrt{(1 + \cos\Phi_{ij})^2\sinh^2 t + 2(1 + \cos\Phi_{ij})},
\]
\[
f_2(t) = \sqrt{(1 + \cos\Phi_{ik})^2\sinh^2 t + 2(1 + \cos\Phi_{ik})}.
\]
By taking the derivative,
\begin{equation}\label{eq:5.2}
f'(t) = \frac{a(\Phi)\sinh(2t)[b(\Phi)\sinh^2 t + c(\Phi)]}{f_1^3(t)f_2^3(t)},
\end{equation}
where
\begin{align*}
a(\Phi) &= (1 + \cos\Phi_{ij})(1 + \cos\Phi_{ik}), \\
b(\Phi) &= (1 + \cos\Phi_{ij})(1 + \cos\Phi_{ik})(1 + \cos\Phi_{jk}), \\
c(\Phi) &= (1 + \cos\Phi_{jk})(2 + \cos\Phi_{ij} + \cos\Phi_{ik}) - (\cos\Phi_{ij} - \cos\Phi_{ik})^2.
\end{align*}
Note the weight function $\Phi : E \to [0,\pi/2]$, hence $\cos\Phi_{ij}, \cos\Phi_{ik}, \cos\Phi_{jk} \in [0,1]$. This implies that
$f_1(t) > 0, f_2(t) > 0, a(\Phi) > 0, b(\Phi) > 0$ and
\begin{align*}
c(\Phi) &= (1 + \cos\Phi_{jk})(2 + \cos\Phi_{ij} + \cos\Phi_{ik}) + 2\cos\Phi_{ij}\cos\Phi_{ik} \\
&\quad + \cos\Phi_{ij}(1 - \cos\Phi_{ij}) + \cos\Phi_{ik}(1 - \cos\Phi_{ik}) \geq 2.
\end{align*}
Therefore, $f'(t) > 0$ and $f(t)$ is strictly increasing, which is equivalent to that $\theta_i(t\vec{1})$ is strictly decreasing.
\end{proof}

\begin{prop}\label{prop:5.4}
For any weight function $\Phi: E \to [0,\pi/2]$, the following limits exist:
\[
\lim_{t \to 0} \theta_i(t\vec{1}) = \arccos \frac{1 + \cos\Phi_{ij} + \cos\Phi_{ik} - \cos\Phi_{jk}}{2\sqrt{1 + \cos\Phi_{ij}}\sqrt{1 + \cos\Phi_{ik}}},
\]
where $\vec{1} = (1,1,1)$, and $\lim_{t \to +\infty} \theta_i(t\vec{1}) = 0$. As a special case, if the weight $\Phi: E \to [0,\pi/2]$ is a constant,
we obtain $\lim_{t \to 0} \theta_i(t\vec{1}) = \frac{\pi}{3}$.
\end{prop}

\begin{proof}
Taking limits directly in the formula \eqref{eq:5.1}, we get the first two limits. In the case that $\Phi: E \to
[0,\pi/2]$ is a constant, we set $\Phi \equiv \varphi$, $\forall e_{ij} \in E$. By \eqref{eq:5.1} we have
\[
f(t) = \frac{(1 + \cos\varphi)^2\sinh^2 t + (1 + \cos\varphi)}{(1 + \cos\varphi)^2\sinh^2 t + 2(1 + \cos\varphi)}.
\]
Then the third limit can follow from direct computations.
\end{proof}

Now we are ready to prove the following comparison principle for inner angles which is key to proofs
of our main theorems.

\begin{lem}\label{lem:5.6}
For a fixed hyperbolic triangle $\triangle ijk$. If $r_i = \min\{r_i, r_j, r_k\}$, then $\theta_i(\vec{r}) \geq \theta_i(r_i\vec{1})$, while if
$r_i = \max\{r_i, r_j, r_k\}$, then $\theta_i(\vec{r}) \leq \theta_i(r_i\vec{1})$.
\end{lem}

\begin{proof}
Let $\sigma: [0,1] \to \mathbb{R}^3_{>0}$ be a curve defined as
\[
\sigma(s) = (r_i(s), r_j(s), r_k(s)) := (1-s)\vec{r} + s r_i \vec{1}, \quad \forall s \in [0,1],
\]
which connects the two points $\vec{r}, r_i\vec{1} \in \mathbb{R}^3_{>0}$ and
\[
r_i(s) = r_i, \quad r_j(s) = r_j + s(r_i - r_j), \quad r_k(s) = r_k + s(r_i - r_k), \quad \forall s \in [0,1].
\]
If $r_i = \min\{r_i, r_j, r_k\}$, then by Lemma \ref{lem:Thurston},
\begin{equation}\label{eq:5.3}
\frac{d}{ds}(\theta_i(\sigma(s))) = (r_i - r_j)\frac{\partial\theta_i}{\partial r_j}(\sigma(s)) + (r_i - r_k)\frac{\partial\theta_i}{\partial r_k}(\sigma(s)) < 0.
\end{equation}
Hence $\theta_i(\sigma(s))$ is decreasing in $[0,1]$. This yields that
\[
\theta_i(\vec{r}) = \theta_i(\sigma(0)) \geq \theta_i(\sigma(1)) = \theta_i(r_i\vec{1}).
\]
Similarly, if $r_i = \max\{r_i, r_j, r_k\}$, then by Lemma \ref{lem:Thurston},
\begin{equation}\label{eq:5.4}
\frac{d}{ds}(\theta_i(\sigma(s))) = (r_i - r_j)\frac{\partial\theta_i}{\partial r_j}(\sigma(s)) + (r_i - r_k)\frac{\partial\theta_i}{\partial r_k}(\sigma(s)) > 0.
\end{equation}
Hence $\theta_i(\sigma(s))$ is increasing in $[0,1]$. This yields that
\[
\theta_i(\vec{r}) = \theta_i(\sigma(0)) \leq \theta_i(\sigma(1)) = \theta_i(r_i\vec{1}).
\]
This completes the proof.
\end{proof}

\section{Long time existence}\label{sec:longtime}

Chow-Luo \cite{ChowLuo} proved the long time existence of the solutions to the combinatorial Ricci flow \eqref{eq:hyperbolic-flow} by
the maximum principle. In this section we give an alternative proof based on our comparison principle. We
need the following calculus lemma introduced by Ge-Hua \cite{GeHua}. For a continuous function $f: [0,+\infty) \to \mathbb{R}$
and any $C \in \mathbb{R}$, the upper level set of $f$ at $C$ is defined as $\{f > C\} := \{t \in [0,+\infty) : f(t) > C\}$. The
lower level set $\{f < C\}$ is defined similarly.

\begin{lem}[Lemma~3.9, \cite{GeHua}]\label{lem:6.1}
Let $f: [0,+\infty) \to \mathbb{R}$ be a locally Lipschitz function. Suppose that
there is a constant $C_1$ such that $f'(t) \leq 0$, for a.e. $t \in \{f > C_1\}$, then
\[
f(t) \leq \max\{f(0), C_1\}, \quad \forall t \in [0,+\infty).
\]
Similarly, if $f'(t) \geq 0$, for a.e. $t \in \{f < C_1\}$, then
\[
f(t) \geq \min\{f(0), C_1\}, \quad \forall t \in [0,+\infty).
\]
\end{lem}

To estimate the inner angles in hyperbolic geometry, we need the following proposition in \cite{ChowLuo} on a
hyperbolic triangle $\triangle ijk$ in $\mathbb{H}^2$.

\begin{lem}[Lemma~3.5, \cite{ChowLuo}]\label{lem:6.2}
For any $\epsilon > 0$, there exists a constant $C_2 = C_2(\epsilon)$ such that when
$r_i > C_2$, the inner angle $\theta_i$ in the hyperbolic triangle $\triangle ijk$ is smaller than $\epsilon$.
\end{lem}

Chow-Luo \cite{ChowLuo} once obtained the following theorem. We give an alternative proof here.

\begin{prop}[Corollary~3.6, \cite{ChowLuo}]\label{prop:6.3}
Let $r(t)$ be a solution to the combinatorial Ricci flow \eqref{eq:hyperbolic-flow}. Then
there exists a positive constant $C_3 = C_3(\mathcal{T}, r(0)) > 0$ depending on the triangulation $\mathcal{T}$ and the initial
data $r(0)$ such that $r_i(t) \leq C_3$ for all $i \in V$.
\end{prop}

\begin{proof}
Set $f(t) := \max_{m \in V} r_m(t)$. Then $f(t)$ is a locally Lipschitz function and for a.e. $t \in [0,+\infty)$,
there exists $i \in V$ depending on $t$, such that
\begin{equation}\label{eq:6.1}
f(t) = r_i(t), \quad f'(t) = r_i'(t).
\end{equation}
Let $C_2$ be the constant determined in Lemma \ref{lem:6.2} such that for any hyperbolic triangle $\triangle ijk$, if $r_i \geq C_2$,
then $\theta_i \leq \pi / \max_{m \in V} d_m$. We claim that $f'(t) \leq 0$, for a.e. $t \in \{f > C_2\}$.

Let $t \in [0,+\infty)$ and $i \in V$ satisfying \eqref{eq:6.1} and $t \in \{f > C_2\}$. Then for any hyperbolic triangle $\triangle ijk$
incident to $i$ realized by the circle packing $r(t)$, $\theta_i(r(t)) \leq \pi / \max_{m \in V} d_m$. Hence by the definition of the
combinatorial Gauss curvature,
\[
K_i(r(t)) = 2\pi - \sum_{\triangle ijk \in F} \theta_i(r(t)) > \pi.
\]
Then by the combinatorial Ricci flow \eqref{eq:hyperbolic-flow}, $f'(t) = r_i'(t) = -K_i \sinh r_i < 0$. This proves the claim.
Then the theorem follows from Lemma \ref{lem:6.1}.
\end{proof}

We provide an alternative proof for the long time existence of the solution to \eqref{eq:hyperbolic-flow}.

\begin{prop}[Proposition~3.4, \cite{ChowLuo}]\label{prop:6.4}
For any initial circle packing $r(0) \in \mathbb{R}^N_{>0}$, the solution $r(t)$ to
the combinatorial Ricci flow \eqref{eq:hyperbolic-flow} exists for all time $t \in [0,+\infty)$.
\end{prop}

\begin{proof}
Since all functions in the equation \eqref{eq:hyperbolic-flow} are smooth and locally Lipschitz continuous, there is
a unique solution $r(t)$ with $t \in [0, T)$ and $0 < T \leq +\infty$ by the classical ODE theory. Rewrite \eqref{eq:hyperbolic-flow} as
$d \ln(\tanh \frac{r_i}{2}) / dt = -K_i$. Note $|K_i| \leq 2\pi \max_{m \in V}(d_m + 1) =: C$. Thus
\[
\tanh\frac{r_i(0)}{2} e^{-Ct} \leq \tanh\frac{r_i(t)}{2} \leq \tanh\frac{r_i(0)}{2} e^{Ct},
\]
which implies that $r_i(t)$ cannot go to $0$ in finite time. On the other hand, by Proposition \ref{prop:6.3}, there exists
a constant $C_3$ depending on the initial data $r(0)$ and the triangulation $\mathcal{T}$ such that $r_i(t) \leq C_3$. Therefore,
by the extension theorem of solutions in ODE theory, the solution to the combinatorial Ricci flow \eqref{eq:hyperbolic-flow}
exists for all $t \in [0,+\infty)$.
\end{proof}

\section{Proof of Theorem \ref{thm:4.1}}\label{sec:proof}

In this section we give the proof of the main theorem, namely Theorem \ref{thm:4.1}. First, we prove the lower
bound estimate for the solution to the combinatorial Ricci flow \eqref{eq:hyperbolic-flow}.

\begin{theo}\label{thm:7.1}
Let $X$ be a closed surface with given weighted triangulation $(\mathcal{T},\Phi)$. Assume the
character $L_i(\mathcal{T},\Phi) > 2\pi$ for all $i \in V$. Let $r(t)$ be a solution to the combinatorial Ricci flow \eqref{eq:hyperbolic-flow}.
Then there is a positive constant $C = C(\mathcal{T}, \Phi, r(0)) > 0$ depending only on the weighted triangulation
$(\mathcal{T},\Phi)$ and the initial data $r(0)$ such that $r_i(t) \geq C$ for all vertices and all times.
\end{theo}

\begin{proof}
Set $g(t) := \min_{m \in V} r_m(t)$. Then $g(t)$ is a locally Lipschitz function and for a.e. $t \in [0,+\infty)$,
there exists a special vertex $i \in V$ depending on $t$, such that
\begin{equation}\label{eq:7.1}
g(t) = r_i(t), \quad g'(t) = r_i'(t).
\end{equation}
For the particular vertex $i$, by Proposition \ref{prop:5.2} and Proposition \ref{prop:5.4}, $\theta_i(t\vec{1})$ is continuously differentiable
and
\[
\lim_{t \to 0} \theta_i(t\vec{1}) = \arccos \frac{1 + \cos\Phi_{ij} + \cos\Phi_{ik} - \cos\Phi_{jk}}{2\sqrt{1 + \cos\Phi_{ij}}\sqrt{1 + \cos\Phi_{ik}}}.
\]
Hence, for any positive constant $\epsilon_0$, small enough and is to be determined later, there exists a constant
$C = C(\mathcal{T}, \Phi, r(0)) > 0$ such that for any $t \leq C$,
\[
\theta_i(t\vec{1}) \geq \arccos \frac{1 + \cos\Phi_{ij} + \cos\Phi_{ik} - \cos\Phi_{jk}}{2\sqrt{1 + \cos\Phi_{ij}}\sqrt{1 + \cos\Phi_{ik}}} - \epsilon_0.
\]
We claim that $g'(t) \geq 0$, for a.e. $t \in \{g < C\}$.

Let $t \in [0,+\infty)$ and $i \in V$ satisfying \eqref{eq:7.1} and $t \in \{g < C\}$. Then for any hyperbolic triangle $\triangle ijk$
incident to $i$ realized by the circle packing $r(t)$,
\[
r_i(t) = \min\{r_i(t), r_j(t), r_k(t)\} < C.
\]
Set $\vec{r}(t) = (r_i(t), r_j(t), r_k(t))$, then by Lemma \ref{lem:5.6},
\[
\theta_i(\vec{r}(t)) \geq \theta_i(r_i(t)\vec{1}) \geq \arccos \frac{1 + \cos\Phi_{ij} + \cos\Phi_{ik} - \cos\Phi_{jk}}{2\sqrt{1 + \cos\Phi_{ij}}\sqrt{1 + \cos\Phi_{ik}}} - \epsilon_0.
\]
Under the assumption $L_j(\mathcal{T},\Phi) > 2\pi$ for all $j \in V$, we choose the constant $\epsilon_0$ so that
\[
0 < \epsilon_0 < \frac{\min_{j \in V}(L_j(\mathcal{T},\Phi) - 2\pi)}{\max_{j \in V} d_j}.
\]
Note $\epsilon_0$ depends on the data of the weighted triangulation $(\mathcal{T},\Phi)$ on $X$. It follows
\begin{align*}
K_i(r(t)) &= 2\pi - \sum_{\triangle ijk \in F} \theta_i(\vec{r}(t)) \\
&\leq 2\pi - \sum_{\triangle ijk \in F} \left(\arccos \frac{1 + \cos\Phi_{ij} + \cos\Phi_{ik} - \cos\Phi_{jk}}{2\sqrt{1 + \cos\Phi_{ij}}\sqrt{1 + \cos\Phi_{ik}}} - \epsilon_0\right) \\
&= \epsilon_0 \cdot d_i - (L_i(\mathcal{T},\Phi) - 2\pi) \\
&\leq \epsilon_0 \cdot \max_{j \in V} d_j - \min_{j \in V}(L_j(\mathcal{T},\Phi) - 2\pi) \\
&< 0.
\end{align*}
By \eqref{eq:7.1} and the combinatorial Ricci flow \eqref{eq:hyperbolic-flow}, $g'(t) = r_i'(t) = -K_i \sinh r_i \geq 0$. This proves the claim.
Then the theorem follows from the claim and Lemma \ref{lem:6.1}.
\end{proof}

Let $r(t)$ be a solution to the combinatorial Ricci flow \eqref{eq:hyperbolic-flow}. Then by Theorem \ref{thm:7.1} and Proposition \ref{prop:6.3},
there exist positive constants $C_1$ and $C_2$ depending on the weighted triangulation $(\mathcal{T},\Phi)$ and the initial
data $r(0)$ such that
\[
C_1 \leq r_i(t) \leq C_2, \quad \forall i \in V, \; t \in [0,+\infty).
\]
This is equivalent to say that $r(t)$ lies in the compact region in $\mathbb{R}^N_{>0}$. Then the existence part of Theorem
\ref{thm:4.1} is obtained by the following proposition proved by Chow-Luo \cite{ChowLuo}.

\begin{prop}[Proposition~3.7, \cite{ChowLuo}]\label{prop:7.2}
Suppose $r(t)$ for $t \in [0,+\infty)$ is a solution to the combinatorial
Ricci flow \eqref{eq:hyperbolic-flow} in the hyperbolic background geometry so that the set $\{r(t) \mid t \in [0,+\infty)\}$ lies in a compact
region in $\mathbb{R}^N_{>0}$. Then there exists a hyperbolic circle packing $r_{ze}$ in $\mathbb{R}^N_{>0}$ and $r(t)$ converges exponentially
fast to $r_{ze}$.
\end{prop}

Obviously, the uniqueness part of Theorem \ref{thm:4.1} is a consequence of the rigidity part of Koebe-Andreev-
Thurston theorem, i.e. the curvature map $r \mapsto K$ is injective. Finally, we prove the nonexistence part of
Theorem \ref{thm:4.1}.

\begin{proof}[Proof of nonexistence part of Theorem \ref{thm:4.1}]
Assume $L_i(\mathcal{T},\Phi) \leq 2\pi$ at each vertex $i \in V$. Hence the average character $L_{av} \leq 2\pi$. By
Proposition \ref{prop:3.4} we get $\chi(X) \geq 0$. For any hyperbolic triangle $\triangle ijk$, denote $\mathrm{Area}(\triangle ijk)$ by its hyperbolic
area. Then $\mathrm{Area}(X) = \sum_{\triangle ijk \in F} \mathrm{Area}(\triangle ijk)$. The combinatorial Gauss-Bonnet formula \cite{ChowLuo} in the
hyperbolic background says:
\[
\sum_{i \in V} K_i = 2\pi\chi(X) + \mathrm{Area}(X).
\]
Hence there is no circle packing $r$ with curvature zero. By Chow-Luo's result \cite{ChowLuo} that the combinatorial
Ricci flow \eqref{eq:hyperbolic-flow} converges if and only if there exists a circle packing $r_{ze}$ with curvature zero, we get the
nonexistence part of Theorem \ref{thm:4.1}.
\end{proof}

If $L_i(\mathcal{T},\Phi) < 2\pi$ for all vertices, we can even show more, that is, any initial circle packing shrinks to
a point along the combinatorial Ricci flow \eqref{eq:hyperbolic-flow}. Actually, we have

\begin{theo}\label{thm:7.3}
If the character $L_i(\mathcal{T},\Phi) < 2\pi$ or $L_i(\mathcal{T},\Phi) \leq 2\pi$ for a constant weight at each vertex,
then the solution $r(t)$ to the combinatorial Ricci flow \eqref{eq:hyperbolic-flow} satisfies $r(t) \to 0$ when $t$ tends to $+\infty$.
\end{theo}

\begin{proof}
We claim that if the character $L_i(\mathcal{T},\Phi) < 2\pi$ or $L_i(\mathcal{T},\Phi) \leq 2\pi$ for a constant weight at each
vertex, then for any circle packing $r$ there exists a vertex $i$ such that
\[
K_i \geq C_0,
\]
where $C_0 = C_0(\mathcal{T},\Phi) > 0$ is a sufficiently small positive constant depending on the weighted triangulation
$(\mathcal{T},\Phi)$.

Let $i$ be the vertex such that $r_i = \max_{j \in V} r_j$. For any hyperbolic triangle $\triangle ijk$ incident to $i$, we have
$r_i = \max\{r_i, r_j, r_k\}$. Set $\vec{r} = (r_i, r_j, r_k)$, by Lemma \ref{lem:5.6},
\begin{equation}\label{eq:7.2}
\theta_i(\vec{r}) \leq \theta_i(r_i\vec{1}).
\end{equation}
If the character $L_j(\mathcal{T},\Phi) < 2\pi$ for each vertex $j \in V$, we choose the constant $\epsilon_0$ depends only on the
data of the weighted triangulation $(\mathcal{T},\Phi)$ on $X$ so that
\[
0 < \epsilon_0 < \frac{\min_{j \in V}(2\pi - L_j(\mathcal{T},\Phi))}{\max_{j \in V} d_j}.
\]
Set
\[
C_0 = \min_{j \in V}(2\pi - L_j(\mathcal{T},\Phi)) - \epsilon_0 \cdot \max_{j \in V} d_j > 0.
\]
By Proposition \ref{prop:5.2} and Proposition \ref{prop:5.4}, $\theta_i(t\vec{1})$ is continuously differentiable and
\[
\lim_{t \to 0} \theta_i(t\vec{1}) = \arccos \frac{1 + \cos\Phi_{ij} + \cos\Phi_{ik} - \cos\Phi_{jk}}{2\sqrt{1 + \cos\Phi_{ij}}\sqrt{1 + \cos\Phi_{ik}}}.
\]
Therefore, for the positive constant $\epsilon_0$, there exists a constant $C > 0$ such that for any $t \leq C$,
\[
\theta_i(t\vec{1}) \leq \arccos \frac{1 + \cos\Phi_{ij} + \cos\Phi_{ik} - \cos\Phi_{jk}}{2\sqrt{1 + \cos\Phi_{ij}}\sqrt{1 + \cos\Phi_{ik}}} + \epsilon_0.
\]
Choose any $t$ satisfying that $0 < t < \min\{r_i, C\}$. Combining \eqref{eq:7.2} and Proposition \ref{prop:5.2} which states
that $t \mapsto \theta_i(t\vec{1})$ is strictly decreasing we have
\begin{equation}\label{eq:7.3}
\theta_i(\vec{r}) \leq \theta_i(r_i\vec{1}) < \theta_i(t\vec{1}) \leq \arccos \frac{1 + \cos\Phi_{ij} + \cos\Phi_{ik} - \cos\Phi_{jk}}{2\sqrt{1 + \cos\Phi_{ij}}\sqrt{1 + \cos\Phi_{ik}}} + \epsilon_0.
\end{equation}
It follows
\begin{align*}
K_i(r) &= 2\pi - \sum_{\triangle ijk \in F} \theta_i(\vec{r}) \\
&\geq 2\pi - \sum_{\triangle ijk \in F} \left(\arccos \frac{1 + \cos\Phi_{ij} + \cos\Phi_{ik} - \cos\Phi_{jk}}{2\sqrt{1 + \cos\Phi_{ij}}\sqrt{1 + \cos\Phi_{ik}}} + \epsilon_0\right) \\
&= (2\pi - L_i(\mathcal{T},\Phi)) - \epsilon_0 \cdot d_i \\
&\geq C_0.
\end{align*}
On the other hand, if the weight function $\Phi : E \to [0,\pi]$ is a constant, then $L_i(\mathcal{T},\Phi) = \frac{\pi d_i}{3}$, so the
assumption $L_i(\mathcal{T},\Phi) \leq 2\pi$ is equivalent to $d_i \leq 6$ at each vertex. By Proposition \ref{prop:5.2} and Proposition \ref{prop:5.4}
we know that the angle function $t \mapsto \theta_i(t\vec{1})$ is continuously differentiable and
\[
\lim_{t \to 0} \theta_i(t\vec{1}) = \frac{\pi}{3}, \quad \lim_{t \to +\infty} \theta_i(t\vec{1}) = 0.
\]
Suppose that $C_0 > 0$ is any fixed sufficiently small positive number and set $\delta_0 = \frac{C_0}{6}$, then by the
intermediate value theorem and above limits, there exists $t_0$, $0 < t_0 < r_i$ such that
\[
\theta_i(t_0\vec{1}) = \frac{\pi}{3} - \delta_0 > 0.
\]
Note that $t \mapsto \theta_i(t\vec{1})$ is strictly decreasing, by \eqref{eq:7.2} we have
\[
\theta_i(\vec{r}) \leq \theta_i(r_i\vec{1}) < \theta_i(t_0\vec{1}) = \frac{\pi}{3} - \delta_0.
\]
Since $d_i \leq 6$,
\[
K_i = 2\pi - \sum_{\triangle ijk \in F} \theta_i(\vec{r}) \geq 2\pi - 6\left(\frac{\pi}{3} - \delta_0\right) = C_0.
\]
This completes the proof of the claim. Let $r(t)$ be a solution to the combinatorial Ricci flow \eqref{eq:hyperbolic-flow}.
Set $r_M(t) := \max_{i \in V} r_i(t)$. Then for a.e. $t \in [0,+\infty)$, there exists $i \in V$ depending on $t$, such that
$r_M(t) = r_i(t)$ and $r_M'(t) = r_i'(t)$. By the claim above, for the circle packing $r(t)$, we have $K_i(r(t)) \geq C_0$,
where $C_0 > 0$ is a sufficiently small positive constant depending on $(\mathcal{T},\Phi)$. By the combinatorial Ricci
flow \eqref{eq:hyperbolic-flow}, for a.e. $t \in [0,+\infty)$
\[
r_M'(t) = r_i'(t) \leq -C_0 \sinh(r_i(t)) = -C_0 \sinh(r_M(t)),
\]
which is equivalent to $(\ln \tanh \frac{r_M(t)}{2})' \leq -C_0$. Integrate both sides from $0$ to $t$, $\tanh \frac{r_M(t)}{2} \leq
\tanh \frac{r_M(0)}{2} e^{-C_0 t}$. Hence $r(t) \to 0$ as $t \to \infty$. This completes the proof of Theorem \ref{thm:7.3}.
\end{proof}

\section{The prescribed flow}\label{sec:prescribed}

In this final section, we derive Theorem \ref{thm:2.7}. We call a given $\bar{K} = (\bar{K}_1, \cdots, \bar{K}_N)$ a prescribed combinatorial
Gaussian curvature. We want to know if there is a circle packing $\bar{r}$ with curvature $K(\bar{r}) = \bar{K}$. This is
called a prescribed circle packing problem. By Thurston's rigidity theorem, the prescribed circle packing
$\bar{r}$ is unique if it exists. By this terminology, a hyperbolic circle packing is just a packing with prescribed
curvature zero.

Obviously, if the prescribed circle packing problem has a solution, then the prescribed curvature $\bar{K}$
lies in $K(\mathbb{R}^N_{>0})$, which is the image set of the curvature map $K$ and vice versa. It has been proved \cite{ChowLuo}
that the prescribed circle packing problem has a solution, i.e. there is a circle packing $\bar{r}$ with curvature
$K(\bar{r}) = \bar{K}$, if and only if the following prescribed combinatorial Ricci flow converges.
\begin{equation}\label{eq:prescribed-flow}
\frac{dr_i}{dt} = (\bar{K}_i - K_i) \sinh r_i.
\end{equation}
From the famous Koebe-Andreev-Thurston theorem, the image set $K(\mathbb{R}^N_{>0})$ of the curvature map
$r \mapsto K$ is a bounded convex polytope. Using the data of $(\mathcal{T},\Phi)$ on $X$, it is described by a set of linear
inequalities:
\begin{equation}\label{eq:image-set}
K(\mathbb{R}^N_{>0}) = \left\{x \in \mathbb{R}^N : \sum_{i \in I} x_i > -\sum_{(e,v) \in \mathrm{Lk}(I)} (\pi - \Phi(e)) + 2\pi\chi(F_I)\right\},
\end{equation}
where $I$ is taken over all proper subsets of $V$, $F_I$ is the subcomplex whose vertices are in $I$, and $\mathrm{Lk}(I)$ is
the set of pairs $(e,v)$ of an edge $e$ and a vertex $v$ satisfying the following: (i) The end points of $e$ are not
in $I$; (ii) $v$ is in $I$; (iii) $e$ and $v$ form a triangle.

\eqref{eq:image-set} describes the image set of $K$ completely. Despite its accuracy, it is global and not easy to verify.
In this section, we will investigate the prescribed flow \eqref{eq:prescribed-flow} further, and prove the following

\begin{theo}\label{thm:8.1}
Consider a given weighted triangulation $(\mathcal{T},\Phi)$ on a closed surface $X$. If
\[
2\pi - L_i(\mathcal{T},\Phi) < \bar{K}_i < 2\pi
\]
at each vertex $i \in V$, then the solution to the prescribed flow \eqref{eq:prescribed-flow} converges exponentially fast to a
packing $\bar{r}$ with curvature $K(\bar{r}) = \bar{K}$. Thus $\bar{r}$ is the solution of the prescribed circle packing problem.

On the other hand, if $\bar{K}_i \leq 2\pi - L_i(\mathcal{T},\Phi)$ at each vertex $i \in V$, then the prescribed flow \eqref{eq:prescribed-flow} can not
converge. Hence the prescribed circle packing problem had no solutions and equivalently, there is no circle
packing $\bar{r}$ with curvature $K(\bar{r}) = \bar{K}$.
\end{theo}

Combining Theorem \ref{thm:8.1} and our ``degree-type'' criteria, i.e. Corollary \ref{cor:3.5}, \ref{cor:3.7} and \ref{cor:3.8}, we have some
quite interesting results for the image set.

\begin{cor}\label{cor:8.3}
Assume the weight function $\Phi: E \to [0,\pi/2]$ is a constant, then
\[
\prod_{i \in V}\left(2\pi - \frac{d_i \pi}{3}, 2\pi\right) \subset K(\mathbb{R}^N_{>0})
\]
and
\[
\prod_{i \in V}\left(-\infty, 2\pi - \frac{d_i \pi}{3}\right] \cap K(\mathbb{R}^N_{>0}) = \emptyset.
\]
\end{cor}

\begin{cor}\label{cor:8.4}
Assume the weight $\Phi \in [0,\pi/2]$ and $d_i \geq 9$ for each vertex $i \in V$, then
\[
[-0.07\pi, 2\pi)^N \subset \left(2\pi - 9\arccos\frac{3}{4}, 2\pi\right)^N \subset K(\mathbb{R}^N_{>0}).
\]
\end{cor}

\begin{cor}\label{cor:8.5}
Assume $d_i \geq 7$ for each vertex $i \in V$. If further assume the weight function $\Phi :
E \to [0,\pi/2]$ satisfies: $\Phi$ is a constant, or $\Phi \in [0, 0.33\pi]$, or $\Phi \in [0.4\pi, \pi/2]$, or more generally $\Phi \in
[\arccos\eta, \arccos\xi] \subset [0,\pi/2]$, where $0 \leq \xi \leq \eta \leq 1$ are arbitrary chosen so that $\eta < (2\cos\frac{2\pi}{7} - 1 + \xi)/(2 -
2\cos\frac{2\pi}{7})$. Then $[0, 2\pi)^N \subset K(\mathbb{R}^N_{>0})$.
\end{cor}

\section*{Acknowledgements}

We are very grateful to the referees for carefully reading the original manuscript and
pointing out some typos. H. Ge is supported by National Natural Science Foundation of China (Grant No.
11871094 and Grant No. 12122119). A. Lin is supported by National Natural Science Foundation of China
(Grant No. 12171480), Hunan Provincial Natural Science Foundation of China (Grant No. 2020JJ4658 and
Grant No. 2022JJ10059) and Scientific Research Program Funds of NUDT (Grant No. 22-ZZCX-016). The
second author would like to thank professor K. Feng and Z. Zhou for communications on related topics. Both
authors would like to thank professor L. Shen for many useful conversations.

\begin{thebibliography}{99}

\bibitem{Andreev1} Andreev E. On convex polyhedra in Loba\v{c}evski\u{i} spaces (Russian). Mat Sb (123), 1970, 81: 445--478.

\bibitem{Andreev2} Andreev E. On convex polyhedra of finite volume in Loba\v{c}evski\u{i} spaces (Russian). Mat Sb (125), 1970, 83: 256--260.

\bibitem{BaoBonahon} Bao X, Bonahon F. Hyperideal polyhedra in hyperbolic 3-space. Bull Soc Math France, 2002, 130: 457--491.

\bibitem{BobenkoHoffmannSpringborn} Bobenko A, Hoffmann T, Springborn B. Minimal surfaces from circle patterns: geometry from combinatorics. Ann of Math, 2006, 164: 231--264.

\bibitem{BobenkoSpringborn} Bobenko A, Springborn B. Variational principles for circle patterns and Koebe's theorem. Trans Amer Math Soc, 2004, 356: 659--689.

\bibitem{Bowditch} Bowditch B. Singular Euclidean structures on surfaces. J Lond Math Soc, 1991, 44: 553--565.

\bibitem{BowersStephenson1} Bowers P, Stephenson K. The set of circle packing points in the Teichm\"uller space of a surface of finite conformal type is dense. Math Proc Cambridge Philos Soc (3), 1992, 111: 487--513.

\bibitem{DaiGuLuo} Dai J, Gu X, Luo F. Variational principles for discrete surfaces. Beijing: Higher Education Press, 2008.

\bibitem{ChowLuo} Chow B, Luo F. Combinatorial Ricci flows on surfaces. J Differential Geom (1), 2003, 63: 97--129.

\bibitem{Colin} de Verdi\`ere Y C. Un principe variationnel pour les empilements de cercles. Invent Math (3), 1991, 104: 655--669.

\bibitem{EdmondsEwingKulkarni} Edmonds A, Ewing J, Kulkarni R. Regular tessellations of surfaces and $(p,q,2)$-triangle groups. Ann of Math (1), 1982, 116: 113--132.

\bibitem{FengGeHua} Feng K, Ge H, Hua B. Combinatorial Ricci flows and the hyperbolization of a class of compact 3-manifolds. Geom Topol (3), 2022, 26: 1349--1384.

\bibitem{GeHua} Ge H, Hua B. 3-dimensional combinatorial Yamabe flow in hyperbolic background geometry. Trans Amer Math Soc (7), 2020, 373: 5111--5140.

\bibitem{GeHuaZhou1} Ge H, Hua B, Zhou Z. Circle patterns on surfaces of finite topological type. Amer J Math (5), 2021, 143: 1397--1430.

\bibitem{GeHuaZhou2} Ge H, Hua B, Zhou Z. Combinatorial Ricci flows for ideal circle patterns. Adv Math, 2021, 383: 107698.

\bibitem{GeJiang1} Ge H, Jiang W. On the deformation of inversive distance circle packings II. J Funct Anal (9), 2017, 272: 3573--3595.

\bibitem{GeJiangShen} Ge H, Jiang W, Shen L. On the deformation of ball packings. Adv Math, 2022, 398: 108192.

\bibitem{HeLiu1} He Z, Liu J. On the Teichm\"uller theory of circle patterns. Trans Amer Math Soc (12), 2013, 365: 6517--6541.

\bibitem{HeSchramm1} He Z, Schramm O. Fixed points, Koebe uniformization and circle packings. Ann of Math (2), 1993, 137: 369--406.

\bibitem{HeSchramm2} He Z, Schramm O. Rigidity of circle domains whose boundary has $\sigma$-finite linear measure. Invent Math (2), 1994, 115: 297--310.

\bibitem{HeSchramm3} He Z, Schramm O. Koebe uniformization for ``almost circle domains''. Amer J Math (3), 1995, 117: 653--667.

\bibitem{HuangLiu} Huang X, Liu J. Characterizations of circle patterns and finite convex polyhedra in hyperbolic 3-space. Math Ann, 2017, 368: 213--231.

\bibitem{Koebe} Koebe P. Kontaktprobleme der konformen Abbildung. Ber S\"achs Akad Wiss Leipzig, Math-Phys Kl, 1936, 88: 141--164.

\bibitem{Leibon} Leibon G. Characterizing the Delaunay decompositions of compact hyperbolic surfaces. Geom Topol, 2002, 6: 361--391.

\bibitem{LiuZhou} Liu J, Zhou Z. How many cages midscribe an egg. Invent Math, 2016, 203: 655--673.

\bibitem{Luo} Luo F. A combinatorial curvature flow for compact 3-manifolds with boundary. Electron Res Announc Amer Math Soc, 2005, 11: 12--20.

\bibitem{MardenRodin} Marden A, Rodin B. On Thurston's formulation and proof of Andreev's theorem. In: Computational methods and function theory (Valparaiso, 1989). Lect Notes in Math, vol. 1435. Berlin: Springer, 1990, 103--115.

\bibitem{Martelli} Martelli B. An introduction to geometric topology. \url{http://people.dm.unipi.it/martelli/geometric_topology.html}, 2016.

\bibitem{Rivin1} Rivin I. Euclidean structures on simplicial surfaces and hyperbolic volume. Ann of Math, 1994, 139: 553--580.

\bibitem{Rivin2} Rivin I. A characterization of ideal polyhedra in hyperbolic 3-space. Ann of Math, 1996, 143: 51--70.

\bibitem{RivinHodgson} Rivin I, Hodgson C. A characterization of compact convex polyhedra in hyperbolic 3-space. Invent math, 1993, 111: 77--111.

\bibitem{RodinSullivan} Rodin B, Sullivan D. The convergence of circle packings to the Riemann mapping. J Differential Geom, 1987, 26: 349--360.

\bibitem{Rousset} Rousset M. Sur la rigidit\'e de poly\`edres hyperboliques en dimension 3: cas de volume fini, cas hyperid\'eal, cas fuchsien. Bull Soc Math France (2), 2004, 132: 233--261.

\bibitem{Schlenker} Schlenker J. Hyperideal circle patterns. Math Res Lett (1), 2005, 12: 85--102.

\bibitem{Schramm1} Schramm O. How to cage an egg. Invent Math, 1992, 107: 543--560.

\bibitem{Stephenson} Stephenson K. Introduction to circle packing. The theory of discrete analytic functions. Cambridge: Cambridge University Press, 2005.

\bibitem{Thurston} Thurston W. Geometry and topology of 3-manifolds. Princeton Lecture Notes, 1976.

\bibitem{Thurston2} Thurston W. The finite Riemann mapping theorem. International Symposium at Purdue University on the occasion of the proof of the Bieberbach Conjecture, 1985.

\bibitem{Zhou1} Zhou Z. Circle patterns with obtuse exterior intersection angles. arXiv:1703.01768v3, 2017.

\bibitem{Zhou2} Zhou Z. Producing circle patterns via configurations. arxiv:2010.13076v9, 2021.

\end{thebibliography}

\end{document}
